%% \subsection{Statistical analyses}
\label{lab:DataScrambling_stat}

%% chi2 analysis.

The chi-squared ($\chi^{2}$) fit method is used to determine a signal strengh from a givent
 dataset (a double ratio). This method is commonly employed in various scientific 
 disciplines to assess how well a model describes the observed data.
 The formula for the chi-squared statistic is:
 \begin{equation}
    \chi^2 = \sum \frac{(O_i - E_i)^2}{\sigma_i^2},
 \end{equation}

 where 
 \begin{itemize}
    \item $i$ is the bin index, and it's range is [0-99] becuase we use 100 bins.
    \item $O_i$ is the observed value for data point $i$.
    \item $E_i$ is the expected value for data point $i$.
    \item $\sigma_i^2$ is the uncertainty (standard deviation) associated with the observed value.
 \end{itemize}

Datasets and theory models:
\begin{itemize}
    \item Dataset is a  double ratio obserable binned in a probing perioud.
     Total number of bin is 100 (see Section.~\ref{lab:ToyStudy}). The dataset is eigher toy, scrambled data or real data when
     this analysis is unblinded.
    
     \item Thery models are analytic fuctions predicted by SME theory described in 
     Section~\ref{lab:Theory}. Each model is fitted to a dataset independent to other models.
     and Each model has a free parameter, $\mu$, which is the signal strength.
    
    \item Degree of freedom: 100 bins minus number of free parameters which one in our case.
     So the DoF is 99.
\end{itemize}
